% Generated draft: requires mathematical and source review.
\documentclass[11pt]{article}
\usepackage[margin=25mm]{geometry}
\usepackage{fontspec}
\setmainfont{Latin Modern Roman}
\usepackage{amsmath,amssymb,mathtools,unicode-math}
\setmathfont{Latin Modern Math}
\usepackage{enumitem,needspace,xurl,hyperref}
\setlength{\parindent}{0pt}
\setlength{\parskip}{5pt}
\setlength{\emergencystretch}{3em}
\newcommand{\N}{\mathbb{N}}
\newcommand{\Z}{\mathbb{Z}}
\newcommand{\Q}{\mathbb{Q}}
\newcommand{\R}{\mathbb{R}}
\newcommand{\C}{\mathbb{C}}
\newcommand{\F}{\mathbb{F}}
\newcommand{\A}{\mathbb{A}}
\newcommand{\PP}{\mathbb{P}}
\newcommand{\E}{\mathbb{E}}
\newcommand{\eps}{\varepsilon}
\newcommand{\dd}{\,\mathrm{d}}
\begin{document}
\section*{UnsolvedMath ID: 6700041}
\textbf{Source problem code:} catalogue entry
\subsection{Definitions and mathematical statement}
For a Riemannian metric $g$ on $M$, scalar curvature positivity means $\operatorname{Sc}_g(x)>0$ at every point. Positive Ricci curvature means
\[
\operatorname{Ric}_g(v,v)>0\qquad\text{for every nonzero tangent vector }v.
\]
Since $\operatorname{Sc}_g=\operatorname{tr}_g\operatorname{Ric}_g$, positive Ricci curvature implies positive scalar curvature. It is a stronger hypothesis, not a change of notation.

The source asks whether it would be ``more prudent'' to replace an earlier condition $\operatorname{Sc}(g)>0$ by $\operatorname{Ric}_g>0$. The earlier proposed theorem or construction is absent, so no antecedent or conclusion is specified for that replacement. This is a methodological question about strengthening assumptions, rather than a conjecture that the two curvature conditions are equivalent or that every positive-scalar-curvature manifold admits a positive-Ricci metric.

The summary mentions Bonnet--Myers motivations, but pointwise positive Ricci curvature is not by itself a stated uniform positive lower bound on a noncompact manifold. Any compactness or diameter conclusion needs its own completeness and uniform-lower-bound hypotheses. No such conclusion is added to this fragment.
\subsection{Short English statement}
Would strengthening a preceding scalar-curvature hypothesis to positive Ricci curvature make the intended programme more appropriate? The preceding statement and a criterion for prudence are not supplied.
\subsection{Sources}
\begin{itemize}
\item Local UnsolvedMath registry, permanent ID 6700041, exact source code catalogue entry; proposed by Misha Gromov, 2017. Registry status partially solved, classification PARTIAL-PROGRESS. The source printed tag is [?43], occurrence 43.
\item Misha Gromov, \emph{101 Questions, Problems and Conjectures around Scalar Curvature} (2017), printed tag [?43], PDF page 48, \url{https://www.ihes.fr/~gromov/wp-content/uploads/2018/08/101-problemsOct1-2017.pdf}. Supplied normalized pdf-text extraction accessed July 29, 2026; incomplete author-hosted PDF. No PDF was accessed here.
\item The supplied research summary calls this partial/interpretive, mentions Myers and Bonnet and the stronger nature of positive Ricci curvature, and says prudence is a design choice rather than a theorem. No precise preceding target or reference is supplied.
\item Provenance flags current status and statement-reuse terms NEEDS\_REVIEW; the source marks the passage as a numbered open item.
\end{itemize}
\end{document}
